\begin{textsample}{POS Dim 1 – es – Score 28.00 – es/es\_ef\_anos\_iniciais\_linguagens\_lp\_3.txt}  \label{ex:f1_pos_140}
| Eixo | Descrição | |---|---| | Oralidade | Abrange fala/escuta, atitudes dos interlocutores, compreensão, funcionamento e produção dos gêneros orais, explorados nas suas características discursivas. | | Fala/Escuta | Há atenção às atitudes dos interlocutores, \textbf{imprescindível} para a participação ativa em sociedade.
Nos anos iniciais, são propostas práticas de convivência oral.
Nos anos finais, privilegiam -se as \textbf{exposições} e apresentações, com ênfase nos gêneros orais da língua materna, para atingir identidade social e cultural.
Cumpre \textbf{salientar}, neste âmbito, a importância do trabalho humanizado e democrático com a variação linguística, traço característico de todas as línguas vivas, \textbf{que} deriva de fatores geográficos, socioeconômicos, etários, relativos ao nível de escolaridade, a contextos de fala, dentre outros.
Esse item está presente em todos os anos do Ensino Fundamental, e precisa ser tratado com cautela, valorizando -se a pluralidade linguística presente em todo o território brasileiro, \textbf{marca} distintiva de seus \textbf{sujeitos}.
Deste modo, \textbf{urge}, no contexto escolar, o permanente combate ao preconceito linguístico, \textbf{que}, como qualquer outra manifestação discriminatória, sumariza misérias humanas e sociais.
Este fenômeno fragiliza relações estabelecidas por meio da fala e da escrita, oprime e segrega indivíduos, sobretudo aqueles cujas variantes linguísticas mais se distanciam das prescrições da gramática normativa ( BAGNO, 1999 ), o \textbf{que} contribui para vilipendiar o fundamental direito à linguagem, previsto na Declaração Universal dos Direitos Linguísticos, para a qual toda variante linguística é lógica e inteligível.
Portanto, é necessário \textbf{que} o trabalho com língua portuguesa dentro das escolas se articule às \textbf{teorias} linguísticas \textbf{que} versam sobre a prática de ensino de língua materna e sobre o preconceito linguístico, o \textbf{que}, por si só, não logra êxito, mas, consorciado a outras ações, pode fazer avançar este começo de redefinição de políticas linguísticas, \textbf{que} parece ainda estagnado pelo establishment social \textbf{que} renega as formas linguísticas não validadas e os \textbf{sujeitos} \textbf{que} delas se valem.
Assim, no eixo Fala/Escuta, é impreterível o papel da instituição escolar na mediação dos processos de falar e escutar, conscientizando os \textbf{discentes} para o respeito irrestrito a esses processos e promovendo permanente discussão sobre as muitas relações entre língua e sociedade, com vistas a construir reflexões sobre o \textbf{que} cada \textbf{uma} dessas esferas representa. | | Leitura | Abrange compreensão, interpretação e estratégias de leitura de textos \textbf{literários}, não \textbf{literários} verbais, verbo-visuais, multimodais, midiáticos, em várias esferas de circulação.
O objetivo da leitura é sempre o mesmo : formar \textbf{um} cidadão reflexivo, crítico, ativo, com consciência e autonomia, capaz de pensar e intervir no meio onde \textbf{vive}, transformando a realidade \textbf{que} o cerca.
Ressaltam -se, ainda, neste documento, estratégias de leitura como ativação de conhecimentos prévios, formulação e verificação de hipóteses, compreensão global, localização, retomada de informações, inferências.
A leitura está presente na vida escolar do estudante, mesmo antes de ele \textbf{conseguir} decodificar palavras.
Trata -se de \textbf{um} eixo muito importante, \textbf{visto} \textbf{que} permite ao aluno transitar com fluidez entre as práticas diversas de leitura e escrita, e é instrumento para a construção de conhecimento em todos os componentes curriculares.
Por isso, há \textbf{uma} preocupação com a fluência.
Na essência do eixo Leitura estão os gêneros textuais.
Nos anos iniciais, o trabalho de leitura permite \textbf{que} o estudante se aproprie de \textbf{um} instrumento e, nos anos finais, desenvolva autonomia nos procedimentos e estratégias próprias do ato leitor. | | Eixo | Descrição | |---|---| | Escrita | Compreende a produção de textos não \textbf{literários} verbais, verbo-visuais, multimodais, midiáticos, além de estratégias de planejamento, revisão, reescrita, e avaliação, adequados ao contexto de produção, ao uso da variedade linguística apropriada a esse contexto, aos enunciadores do discurso ( \textbf{autor} e possível leitor ), ao gênero textual, ao suporte e à esfera de circulação.
Abrange ainda a edição de textos nos meios digitais.
A escrita é \textbf{um} processo complexo \textbf{que} \textbf{exige} \textbf{um} projeto de texto organizado, a partir de \textbf{um} gênero e etapas de reescrita, além de distribuição gráfica e \textbf{marcas} de segmentação.
Por isso, este eixo está associado a todos os outros.
Nos dois primeiros anos, temos a importância do professor escriba, \textbf{que} registra a produção individual e coletiva, permitindo \textbf{que} os estudantes tenham a vivência da escrita e a experiência da criação, antes de se apropriarem do sistema de escrita alfabético. | | Análise Linguística/Semiótica | Definimos, conforme as prescrições dos documentos oficiais assumidos neste currículo, o ensino da gramática a partir de reflexões como “ o \textbf{que} ensinar? ”, “ para \textbf{que} ensinar? ”, “ como ensinar? ”, valorizando \textbf{um} trabalho epilinguístico³ nos anos iniciais, a fim de conquistar \textbf{uma} visão metalinguística⁴, nos anos finais do Ensino Fundamental.
O ensino a \textbf{que} nos referimos no eixo Escrita compreende, deste modo, alfabetização, ortografia, morfossintaxe, recursos coesivos, processos de formação de palavras, construção da frase na norma padrão.
Na visão epilinguística, o trabalho está centrado no uso : estudo de ortografia ( contextualizada, ligada à produção de textos ), acentuação, pontuação, concordância, coesão, processos de formação das palavras.
Na visão metalinguística, o trabalho se concentra na análise dos aspectos \textbf{constitutivos} da língua, sua morfologia e sintaxe, bem como o conhecimento da Nomenclatura Gramatical Brasileira, para empregar a metalinguagem adequada. | ³ A atividade epilinguística é o ato de reflexão e operação sobre a linguagem, estimulada e/ou realizada durante o processo de produção de textos.
É \textbf{focada}, também, na compreensão do uso \textbf{que} se faz dos conceitos linguísticos presentes em \textbf{uma} dada situação de comunicação, vislumbrando a possibilidade de aproximação entre essas diversas competências, favorecendo, sobretudo, a autonomia e controle da produção textual.
Dessa forma, a atividade epilinguística é a própria atividade de linguagem.
E essa apenas pode ser estudada a partir das relações léxico-gramaticais da língua.
É \textbf{um} trabalho de construção reflexiva sobre \textbf{um} texto \textbf{que} faz a ponte necessária entre o conhecimento linguístico de \textbf{um} falante e a metalinguística.
É, por assim \textbf{dizer}, \textbf{uma} estratégia \textbf{imprescindível} na ligação entre a capacidade do estudante de produzir textos e a de descrever os fatos linguísticos levados em conta em sua produção.
As atividades epilinguísticas têm como objetivo proporcionar ao usuário da língua oportunidade para refletir sobre os recursos \textbf{expressivos} de \textbf{que} faz uso ao falar ou escrever, tornando consciente a utilização de certos conceitos não didaticamente \textbf{explicitados} a priori, mas aprendidos e ampliados no processo de operacionalização destes no momento da produção de texto.
Esta ferramenta didático-metodológica, denominada epilinguística, contribui significativamente para o ensino de Língua Portuguesa na educação básica, pois \textbf{certamente} possibilita a aliança entre leitura, produção textual e análise linguística. ⁴ A atividade metalinguística deve ser o final do processo, mas não devemos nem podemos estabelecer quando se devem começar tais atividades.
O aluno acostumado a trabalhar com atividades epilinguísticas sentirá, por si só, a necessidade de chegar a conclusões sobre \textbf{uma} \textbf{teoria} gramatical.
Mas não se tem condições de, a priori, determinar quando será esse momento, pois isso \textbf{depende} de certo grau de \textbf{maturidade} linguística.
Sabe -se, portanto, apenas o “ como ”.
Chega -se à metalinguagem “ como resultado de \textbf{uma} familiaridade com os fatos da língua ” e como decorrência de \textbf{uma} necessidade de \textbf{sistematizar} \textbf{um} “ saber linguístico \textbf{que} se \textbf{aprimorou}. ” ( FRANCHI, 1987, p. 41 ). | Eixo | Descrição | |---|---| | Educação \textbf{Literária} | Ressalta -se a importância do texto \textbf{literário}, como objeto artístico, social, histórico e cultural, para se desenvolver o prazer pela leitura.
Consideramos a literatura \textbf{uma} \textbf{forte} aliada da educação estética, pois o trabalho com a literatura na escola permite ao aluno a compreensão da realidade e possibilita a produção de conhecimento por meio da arte da linguagem.
Há, portanto, a necessidade de se fazer o uso adequado do texto \textbf{literário} em sala de aula, respeitando a sua função estética.
Compreendemos ainda \textbf{que} o estudo da literatura é importante em \textbf{uma} perspectiva analítica, discursiva, política e não mercadológica, nem tampouco utilitária, ou seja, \textbf{que} sirva apenas como pretexto para o desenvolvimento de questões multidisciplinares.
Este eixo aborda apreciação, interpretação e produção de textos \textbf{literários} ( narrativos, poéticos e dramáticos ) de \textbf{autores} brasileiros, capixabas, indígenas, latinos, africanos e clássicos, com todas as suas características, relacionando -os aos eixos de oralidade e conhecimentos linguísticos. |

% matched lemmas: aprimorar, autor, certamente, conseguir, constitutivo, depender, discente, dizer, exigir, explicitar, exposição, expressivo, focar, forte, imprescindível, literário, marca, maturidade, que, salientar, sistematizar, sujeito, teoria, um, uma, urgir, ver, viver
\end{textsample}
